\documentclass[10pt,a4paper]{amsart}
\usepackage[margin=19mm]{geometry}
\usepackage{amsmath,amssymb,mathtools}
\usepackage[scr=rsfs,cal=euler]{mathalfa}
\usepackage{xfrac}
\usepackage[unicode,hidelinks,bookmarks=false]{hyperref}
\newcommand{\ie}{i.e.\ }
\newcommand{\cf}{cf.\ }
\newcommand{\resp}{respectively\ }
\newcommand{\Subsection}{\subsection}
\providecommand{\nobreakdash}{\nobreak\hbox{-}}
\setlength{\emergencystretch}{2em}
\multlinegap0pt
\usepackage{amssymb}
\usepackage{mathtools}
\usepackage{float}
\usepackage[normalem]{ulem}
\newtheorem{theorem}{Theorem}
\newcommand{\N}{\mathbb{N}}
\newcommand{\R}{\mathbb{R}}
\newcommand{\exD}{\mathbb{R}^2 \setminus \overline{D}}
\title{Factorization method for a simply supported obstacle from point source measurements via far--field transformation}
\author{Isaac Harris, Andreas Kleefeld}
\date{Source: arXiv:2603.27647v1 (CC BY 4.0)}
\begin{document}
\maketitle

\noindent\emph{Source and licence.} Factorization method for a simply supported obstacle from point source measurements via far--field transformation. Version arXiv:2603.27647v1 (CC BY 4.0), distributed by arXiv under the Creative Commons Attribution 4.0 International licence, http://creativecommons.org/licenses/by/4.0/, as recorded in the arXiv licence metadata for this exact version and shown on the abstract page https://arxiv.org/abs/2603.27647v1. Full licence notice: \texttt{licenses/arxiv-2603.27647-CC-BY-4.0.txt}.\par\smallskip

\noindent\textit{Editorial scope.} Counted unit: the article's theorem ip-uniqueness (source0.tex lines 264--266). The article's complete original proof of it is delivered with it (source0.tex lines 267--282, one \texttt{proof} environment of 16 lines). No mathematics is written, repaired or re-derived: the statement and the proof are the article's own lines, in the article's own notation, and no retained line is substituted or re-flowed.  Retained uncounted as the source-local context the proof needs: lines 217--219; lines 235--237; lines 245--247.  Nothing was omitted from any retained span, so the package carries no omission record.  No retained line contains a citation, so the wrapper carries no bibliography and no bibliography entry.  No retained line contains a figure environment, an image-inclusion command or drawing code, so no image is delivered and the package declares no asset dependency.  Editorial formatting only, in addition: an AMS \texttt{amsart} preamble that restates the article's own declarations and macros used by the retained lines, namely the environments \texttt{theorem} on separate counters, together with the standard packages those lines need.  The retained environments print in this document as Theorem 1 in order of appearance, whereas the article prints the same items with the numbering of its own full document.  The complete native source of the article as distributed in the arXiv e-print for this exact version is bundled unchanged as \texttt{sources/197369/source0.tex}.
\begin{align} \label{nf-data}
\big\{ u^{\text{scat}}( x , y) \quad \text{ and } \quad \Delta_x u^{\text{scat}}( x , y) \quad : \quad \text{ for all $x,y \in \Gamma$} \big\}
\end{align}

\begin{align} \label{vh-bvp} 
\Delta_x u_\textrm{pr} + k^2 u_\textrm{pr} = 0 \quad \text{in } \,  \exD \quad \text{and } \quad u_\textrm{pr} = \frac{1}{2k^2} \Phi_k ( \cdot \, , y)  \quad \text{on } \,  \partial D
\end{align}

\begin{align}\label{recip}
u_\textrm{pr}(x,y) =  u_\textrm{pr}(y,x) \quad \text{ for all $x,y \in \Gamma$}. 
\end{align}

\begin{theorem} \label{ip-uniqueness}
The near--field data defined by equation \eqref{nf-data} uniquely determines the scatterer. 
\end{theorem}

\begin{proof}
To prove the claim we proceed by way of contradiction. To this end, assume there are two scatterers $D^{(1)}$ and $D^{(2)}$ that produce the same near--field data set in \eqref{nf-data}. This would imply that the corresponding $u_\textrm{pr}^{(1)} ( x , y )$ for $D^{(1)}$ and $u_\textrm{pr}^{(2)} ( x , y )$ for $D^{(2)}$ coincide for all $x,y \in \Gamma$. Since the exterior Dirichlet problem for the Helmholtz equation is well--posed this implies that 
$$u_\textrm{pr}^{(1)} ( x , y ) = u_\textrm{pr}^{(2)} ( x , y ) \quad \text{ for all $y\in \Gamma$ and $x \in \R^2 \setminus \overline{(D^{(1)} \cup D^{(2)})}$}$$
by the unique continuation principle. Now, due to the reciprocity identity in \eqref{recip} we can see that 
$$u_\textrm{pr}^{(1)} ( x , y ) = u_\textrm{pr}^{(2)} ( x , y ) \quad \text{ for all $x,y \in \R^2 \setminus \overline{(D^{(1)} \cup D^{(2)})}$}.$$

Without loss of generality, assume that $\overline{D^{(1)}}$ is not contained in $\overline{D^{(2)}}$, then $\partial D^{(1)} \setminus \overline{D^{(2)}}$ is a non--empty set. Therefore, let $y^* \in \partial D^{(1)} \setminus \overline{D^{(2)}}$ and define the sequence 
$$y_n = y^* + \frac{1}{n} \nu_1(y^*)  \quad \text{for $n \in \N$} $$ 
where $\nu_1$ is the outward unit normal vector for $\partial D^{(1)}$. For $n$ sufficiently large we can assume that the sequence $y_n \in \R^2 \setminus \overline{(D^{(1)} \cup D^{(2)})}$. So by the reciprocity identity we have that
$$u_\textrm{pr}^{(1)} ( x , y_n ) = u_\textrm{pr}^{(2)} ( x , y_n ) \quad \text{ for all $x \in \R^2 \setminus \overline{(D^{(1)} \cup D^{(2)})}$}.$$
This gives that by letting $x \to y^*$ we can obtain
$$ u_\textrm{pr}^{(1)} ( y^* , y_n ) = u_\textrm{pr}^{(2)} ( y^* , y_n )  \quad \text{ and by the boundary condition \eqref{vh-bvp}} \quad \frac{1}{2k^2} \Phi_k ( y^* , y_n ) = u_\textrm{pr}^{(2)} ( y^* , y_n )$$
since $y^* \in \partial D^{(1)}$. This gives a contradiction since 
$$ |\Phi_k ( y^* , y_n )| \to \infty  \quad \text{ whereas } \quad |u_\textrm{pr}^{(2)} ( y^* , y_n )| < \infty  \quad \text{ as $n \to \infty$ }$$
by elliptic regularity. This proves the claim. 
\end{proof}

\end{document}
